\documentclass[aspectratio=169,professionalfonts]{beamer}

\usetheme{CityUHK}

\usepackage[T1]{fontenc}
\usepackage[utf8]{inputenc}
\usepackage{lmodern}
\usepackage{mathtools}
\usepackage{bm}
\usepackage{physics}
\usepackage{tikz}
\usepackage{pgfplots}
\usepackage{siunitx}

\pgfplotsset{compat=1.17}

\title[Lantern Dynamics]{A Random Walk Through Lantern Dynamics}
\subtitle{Fictional models, tidy equations, and surprisingly cooperative plots}
\author[Q. Maple et al.]{Quinn Maple \and Riley Chen \and Morgan Vale}
\institute[CityUHK]{Department of Imaginary Systems, City University of Hong Kong}
\date[May 2027]{14 May 2027}
\conference{7th Symposium on Stochastic Lanterns}
\titlenote{Demo deck note: all names, events, and scientific claims are placeholders.}

\begin{document}

{
\usebackgroundtemplate{\CityUHKFullBackground{\CityUHKClosingBg}}
\setbeamertemplate{footline}{}
\begin{frame}[plain]
  \titlepage
\end{frame}
}

\begin{frame}{Outline}
  \tableofcontents
\end{frame}

\section{Setting}

\begin{frame}{Why Lanterns?}{A harmless fictional motivation}
  \begin{columns}[T,totalwidth=\textwidth]
    \column{0.55\textwidth}
    \begin{block}{Toy problem}
      Suppose a lantern drifts through a quiet city square while its brightness reacts to wind, memory, and moonlight.
      We ask whether the observed glow can reveal the hidden drift pattern.
    \end{block}

    \begin{itemize}
      \item The state variable is the latent position $x_t$.
      \item The observable is the brightness signal $b_t$.
      \item The nuisance parameter is the evening breeze $\eta_t$.
    \end{itemize}

    \column{0.40\textwidth}
    \begin{alertblock}{Main message}
      A clean slide should tell the audience what to look at before asking them to interpret details.
    \end{alertblock}

    \begin{exampleblock}{Template purpose}
      This content is intentionally random, so the layout can be reused for any real topic.
    \end{exampleblock}
  \end{columns}
\end{frame}

\begin{frame}{Ingredients}{Minimal objects for the story}
  \begin{enumerate}
    \item A hidden trajectory $x_0,x_1,\ldots,x_T$ moving through a two-dimensional plaza.
    \item A response function $g(x_t)$ that converts location into observed brightness.
    \item A correction term that remembers whether the lantern was bright a moment ago.
    \item A small amount of noise, because perfectly behaved lanterns are suspicious.
  \end{enumerate}

  \vspace{0.3cm}
  \begin{block}{Example sentence}
    ``The lantern forgets slowly, drifts politely, and reports its location through light.''
  \end{block}
\end{frame}

\section{Figure}

\begin{frame}{Figure with Explanation}{Two invented brightness curves}
  \begin{columns}[T,totalwidth=\textwidth]
    \column{0.55\textwidth}
    \centering
    \begin{tikzpicture}
      \begin{axis}[
        width=0.95\linewidth,
        height=4.6cm,
        xmin=0, xmax=8,
        ymin=0, ymax=1.2,
        axis lines=left,
        xlabel={Evening time $t$},
        ylabel={Brightness $b(t)$},
        tick style={CityUGrey},
        every axis x label/.style={at={(ticklabel* cs:1.0)},anchor=west},
        every axis y label/.style={at={(ticklabel* cs:1.0)},anchor=south},
        grid=major,
        grid style={draw=CityUSand!80!black},
        xtick={0,2,4,6,8},
        ytick={0,0.3,0.6,0.9,1.2}
      ]
        \addplot[very thick,CityURed,samples=160,domain=0:8] {0.62 + 0.28*sin(deg(1.25*x))*exp(-0.08*x)};
        \addplot[very thick,CityUGold!85!black,dashed,samples=160,domain=0:8] {0.48 + 0.42*(1-exp(-0.55*x))*exp(-0.025*x)};
        \legend{Observed glow,Smoothed estimate}
      \end{axis}
    \end{tikzpicture}

    \column{0.40\textwidth}
    \begin{block}{Reading the figure}
      \begin{itemize}
        \item The observed glow oscillates early in the evening.
        \item The smoothed estimate rises more calmly.
        \item The two curves agree once the transient behavior fades.
      \end{itemize}
    \end{block}

    \begin{exampleblock}{Reusable pattern}
      Pair the plot with a short reading guide.
    \end{exampleblock}
  \end{columns}
\end{frame}

\section{Math}

\begin{frame}{Pure Math Example}{A compact fictional model}
  \small
  \begin{block}{Latent process}
    For position $x_t \in \mathbb{R}^2$ and noise $\xi_t$,
    \[
      x_{t+1}
      =
      x_t
      + \alpha
      \begin{pmatrix}
        \cos \theta_t \\
        \sin \theta_t
      \end{pmatrix}
      + \sigma \xi_t .
    \]
  \end{block}

  \begin{alertblock}{Brightness functional}
    \[
      b_t
      =
      \exp\left(-\frac{\norm{x_t-c}^2}{2\rho^2}\right)
      + \lambda b_{t-1}
      + \varepsilon_t,
    \]
    Here $c$ is the plaza center and $\lambda$ controls memory.
  \end{alertblock}

  \begin{exampleblock}{Interpretation}
    Nearby lanterns glow more strongly, with a fading trace of the previous moment.
  \end{exampleblock}
\end{frame}

\section{Logic}

\begin{frame}{Derivation Logic}{Memory parameter}
  \footnotesize
  \textbf{Goal.} Estimate the memory coefficient $\lambda$ in
  \[
    b_t = z_t + \lambda b_{t-1} + \varepsilon_t,
    \qquad z_t \text{ known}.
  \]

  \vspace{-0.25cm}
  \begin{align*}
    r_t &= b_t - z_t,\\
    J(\lambda) &= \sum_{t=1}^{T} (r_t-\lambda b_{t-1})^2,\\
    0 &= -2\sum_{t=1}^{T} b_{t-1}(r_t-\lambda b_{t-1}),\\
    \hat{\lambda}
    &= \frac{\sum_{t=1}^{T} b_{t-1}r_t}
            {\sum_{t=1}^{T} b_{t-1}^2}.
  \end{align*}
\end{frame}

\begin{frame}{Derivation Logic}{How to present the chain}
  \begin{columns}[T,totalwidth=\textwidth]
    \column{0.48\textwidth}
    \begin{block}{Make explicit}
      \begin{itemize}
        \item What is observed
        \item What is assumed known
        \item Which quantity is being estimated
      \end{itemize}
    \end{block}

    \column{0.48\textwidth}
    \begin{alertblock}{Keep short}
      \begin{itemize}
        \item Avoid algebra that does not change the idea.
        \item Use one equation per logical step.
        \item Move edge cases to backup slides.
      \end{itemize}
    \end{alertblock}
  \end{columns}

  \vspace{0.15cm}
  \begin{exampleblock}{Speaking note}
    The audience should feel the derivation moving forward, not just expanding sideways.
  \end{exampleblock}
\end{frame}

\section{Wrap-Up}

\begin{frame}{Summary}{What this random demo illustrates}
  \begin{itemize}
    \item Metadata fields for title, subtitle, authors, institute, date, conference, and title note
    \item A title page using the closing-style background
    \item Content pages using the CityUHK content background
    \item Automatic section pages using the section background
    \item Example slides for bullets, figures, equations, and derivations
  \end{itemize}

  \vspace{0.25cm}
  \begin{block}{Next step}
    Replace the lantern story with your actual research, teaching material, or project update.
  \end{block}
\end{frame}

\CityUHKClosingFrame{Thank You}{Questions and comments}

\end{document}
